
The codebase is engineered to promote modularity and scalability, ensuring that individual components can be updated or reused without compromising the integrity of the entire system. A strict separation of concerns has been enforced by organizing the project into distinct directories, effectively decoupling UI elements, business logic, and external service integrations. All core PWA source code is contained within the \texttt{src} directory, structured as follows:

\begin{description}
    \item[Components:] This directory contains both bespoke project components and those derived from the \textbf{shadcn/ui} library. The \texttt{ui} subfolder is reserved for the foundational shadcn library components, while the \texttt{app} subfolder and the root of the directory house custom-built elements.
    \item[Hooks:] This folder contains custom React Hooks, which encapsulate reusable logic to maintain clean components. They are specifically utilized for \textbf{Device Detection} (adapting the UI for mobile vs. desktop) and \textbf{Toast Notifications} (managing system feedback).
    \item[Pages:] This directory follows a file-based routing logic, containing the main views of the application. The \texttt{app} subfolder organizes core views such as the \textbf{Dashboard}, \textbf{ExplorePaths}, and \textbf{RecordTrip}, ensuring a clear mapping between URL paths and UI controllers.
    \item[Stores:] This directory houses global state management logic (e.g., via Zustand). Key files include \texttt{useAuthStore.ts} for authentication states and \texttt{useDebugStore.ts} for development-time performance monitoring.
    \item[Lib:] Dedicated to utility functions and service configurations. It includes \texttt{utils.ts} for shared helpers and specialized modules like \texttt{weatherService.ts}, which abstracts the logic for processing external meteorological data.
\end{description}

\section{Backend Logic and Database Functions}
The backend logic is implemented using Supabase, leveraging PostgreSQL for data storage and spatial operations (PostGIS), and Edge Functions for complex computational logic. Below is a description of the key functions and triggers implemented.

\subsubsection{Edge Functions}
\begin{itemize}
    \item \textbf{func\_path\_search:} 
    This TypeScript-based Edge Function acts as the main entry point for the route search algorithm. It handles the API request containing origin and destination coordinates and invokes the database RPC \texttt{search\_best\_paths\_rpc} to retrieve candidate paths. 
    Crucially, it implements a ranking algorithm that sorts results based on a composite score derived from:
    \begin{itemize}
        \item \textbf{Geometric Quality:} Derived from the spatial match score (exponential decay).
        \item \textbf{Data Freshness:} Penalizing paths that haven't been updated recently (time-decay factor).
        \item \textbf{Path Status:} Weighting paths classified as "optimal" higher than those with "sufficient" or "bad" conditions.
    \end{itemize}
    \textbf{Code:}
    \begin{lstlisting} 
    import { createClient } from "https://esm.sh/@supabase/supabase-js@2";

    const corsHeaders = {
    "Access-Control-Allow-Origin": "*",
    "Access-Control-Allow-Headers": "authorization, x-client-info, apikey, content-type",
    };

    Deno.serve(async (req) => {
    if (req.method === "OPTIONS") {
        return new Response("ok", { headers: corsHeaders });
    }

    try {
        const supabaseClient = createClient(
        Deno.env.get("SUPABASE_URL") ?? "",
        Deno.env.get("SUPABASE_ANON_KEY") ?? "",
        {
            global: {
            headers: { Authorization: req.headers.get("Authorization")! },
            },
        }
        );

        const { origin_lat, origin_lon, dest_lat, dest_lon } = await req.json();

        const { data: candidates, error } = await supabaseClient.rpc(
        "search_best_paths_rpc",
        { origin_lat, origin_lon, dest_lat, dest_lon, buffer_meters: 20000}
        );

        if (error) throw error;

        if (!candidates || candidates.length === 0) {
        return new Response(JSON.stringify([]), { 
            headers: { ...corsHeaders, "Content-Type": "application/json" } 
        });
        }

        const rankedPaths = candidates.map((path: any) => {
        // Geometric Quality: exp(-0.001 * x)
        let geo_match_score = 0;
        if(path.geo_match_score >= path.geo_match_score_rev){
            geo_match_score = path.geo_match_score_rev;
        }else{
            geo_match_score = path.geo_match_score;
        }
        const geometricQuality = Math.exp(-0.001 * geo_match_score);

        // Days since last status update 
        const daysSinceUpdate = (new Date().getTime() - new Date(path.last_update).getTime()) / (1000 * 3600 * 24);
        // Every 30 days the score lost 10% of his value 
        const freshnessFactor = Math.exp(-0.003 * daysSinceUpdate);

        let ratingStatus = 0;
        switch (path.status) {
            case "optimal": ratingStatus = 1; break;
            case "good": ratingStatus = 0.8; break;
            case "sufficient": ratingStatus = 0.6; break;
            default: ratingStatus = 0.3;
        }

        const finalScore = ((geometricQuality * 0.4) + (ratingStatus * freshnessFactor * 0.6))*100;

        return {
            ...path,
            score: finalScore,
        };
        });
        
        rankedPaths.sort((a: any, b: any) => b.score - a.score);

        return new Response(JSON.stringify(rankedPaths.slice(0, 10)), {
        headers: { ...corsHeaders, "Content-Type": "application/json" },
        status: 200,
        });

    } catch (error) {
        return new Response(JSON.stringify({ error: error.message }), {
        headers: { ...corsHeaders, "Content-Type": "application/json" },
        status: 400,
        });
    }
    });
    \end{lstlisting}
\end{itemize}

\subsubsection{Database Functions (RPCs)}
\begin{itemize}
    \item \textbf{check\_path\_similarity:} 
    A PostGIS spatial function designed to prevent the creation of duplicate paths. Before saving a new route, this function compares the geometry of the new activity against existing paths using the \textbf{Fréchet distance} and length tolerance. If a matching path is found (within the specified tolerance meters), the system associates the activity with the existing path instead of creating a new one. \\
    \textbf{Code:}
    \begin{lstlisting}
    DECLARE
        new_geom GEOMETRY;
        new_len FLOAT;
    BEGIN
        new_geom := ST_SetSRID(ST_GeomFromText(new_geom_wkt), 4326);
        new_len := ST_Length(new_geom::geography);

        RETURN QUERY
        SELECT 
            p.id,
            p.name,
            p.description, 
            p.id_owner,    
            p.is_public,
            ST_FrechetDistance(
                ST_Transform(p.geom::geometry, 3857), 
                ST_Transform(new_geom, 3857)
            ) as similarity_score
        FROM 
            paths p
        WHERE 
            p.geom IS NOT NULL
            AND (p.is_public = true OR p.id_owner = auth.uid()) 
            AND ABS(ST_Length(p.geom::geography) - new_len) < (new_len * length_tolerance_pct)
            AND ST_DWithin(p.geom::geography, new_geom::geography, tolerance_meters * 2)
            AND ST_FrechetDistance(ST_Transform(p.geom::geometry, 3857), ST_Transform(new_geom, 3857)) < tolerance_meters
        ORDER BY similarity_score ASC
        LIMIT 5; 
    END;
    \end{lstlisting}

    \item \textbf{get\_path\_community\_usage:} 
    A utility function that calculates the popularity of a specific path. It counts the number of distinct users (excluding the current requester) who have recorded activities on that path, providing a metric for community usage. \\
    \textbf{Code:}
    \begin{lstlisting}
    DECLARE
    usage_count int;
    BEGIN
    IF p_path_id IS NULL THEN
        RETURN 0;
    END IF;

    SELECT count(distinct user_id)
    INTO usage_count
    FROM activities
    WHERE path_id = p_path_id
    AND user_id != auth.uid(); 
    END;
    \end{lstlisting}
    
    \item \textbf{search\_best\_paths\_rpc:} 
    The core spatial query invoked by the Edge Function. It utilizes PostGIS functions (such as \texttt{ST\_DWithin} and \texttt{ST\_Distance}) to efficiently query the \texttt{paths} table. It identifies all public or user-owned paths that spatially connect the requested origin and destination points within a 20km buffer, calculating the geometric match score for both forward and reverse directions. \\
    \textbf{Code:}
    \begin{lstlisting}
    DECLARE
        origin_pt geography;
        dest_pt geography;
    BEGIN
        origin_pt := ST_SetSRID(ST_MakePoint(origin_lon, origin_lat), 4326)::geography;
        dest_pt := ST_SetSRID(ST_MakePoint(dest_lon, dest_lat), 4326)::geography;

        RETURN QUERY
        SELECT 
        p.id, p.name, p.status, p.last_update,
        (
            ST_Distance(ST_StartPoint(p.geom)::geography, origin_pt) +
            ST_Distance(ST_EndPoint(p.geom)::geography, dest_pt)
        )::float8 as geo_match_score,
        (
            ST_Distance(ST_StartPoint(p.geom)::geography, dest_pt) +
            ST_Distance(ST_EndPoint(p.geom)::geography, origin_pt)
        )::float8 as geo_match_score_rev
        FROM public.paths p
        WHERE 
        (
            p.is_public IS TRUE OR p.id_owner = auth.uid()
        )
        AND
        (
            (
            ST_DWithin(ST_StartPoint(p.geom)::geography, origin_pt, buffer_meters)
            AND
            ST_DWithin(ST_EndPoint(p.geom)::geography, dest_pt, buffer_meters)
            )
            OR 
            (
            ST_DWithin(ST_StartPoint(p.geom)::geography, dest_pt, buffer_meters)
            AND
            ST_DWithin(ST_EndPoint(p.geom)::geography, origin_pt, buffer_meters)
            )
        );
    END;
    \end{lstlisting}
\end{itemize}

\subsubsection{Database Triggers}
\begin{itemize}
    \item \textbf{before\_insert\_set\_last\_update:} 
    Ensures data consistency by automatically initializing the \texttt{last\_update} timestamp to the record's creation time (\texttt{created\_at}) upon insertion. \\
    \textbf{Code:}
    \begin{lstlisting}
    BEGIN
    NEW.last_update := NEW.created_at;
    RETURN NEW;
    END;
    \end{lstlisting}

    \item \textbf{set\_last\_update\_after\_status\_update:} 
    Monitors changes to the \texttt{status} column of a path. Whenever the status is modified (e.g., from "good" to "bad"), this trigger automatically refreshes the \texttt{last\_update} timestamp to \texttt{now()}, ensuring the "Data Freshness" score in the search algorithm remains accurate. \\
    \textbf{Code:}
    \begin{lstlisting}
    BEGIN
    IF OLD.status IS DISTINCT FROM NEW.status THEN
        UPDATE public.paths
        SET last_update = now()
        WHERE id = NEW.id;
    END IF;
    RETURN NEW;
    END;
    \end{lstlisting}
\end{itemize}